% Generated from locale/tr/content/computability/computability-theory/reducibility.tex; do not hand-edit.


\olfileid[tr]{cmp}{thy}{red}
\olsection{İndirgenebilirlik}

\begin{explain}
Artık hesaplanabilir biçimde numaralandırılabilir, fakat hesaplanabilir
olmayan en az bir $K_0$ kümesinin bulunduğunu biliyoruz. Başkalarının da
bulunduğu açık olmalıdır. İndirgenebilirlik yöntemi, sürekli ilk ilkelere
dönmeden başka kümelerin de bu özellikleri taşıdığını göstermenin güçlü bir
yolunu sağlar.

Genel olarak bir $A$ kümesinin $B$ kümesine ``indirgenmesi'', !!{element}s
öğelerinin $B$'de bulunup bulunmadığına verilen yanıtları aynı !!{element}s
öğelerinin $A$'da bulunup bulunmadığına ilişkin yanıtlara dönüştürme
yöntemidir. ``Çok-bir indirgenebilirlik'' denen bir kavrama odaklanacağız;
fakat farklı özelliklere sahip başka birçok indirgenebilirlik kavramı vardır.
İndirgenebilirlik kavramları, verimlilik sorunlarının da hesaba katılması
gereken hesaplama karmaşıklığı incelemelerinde merkezîdir. Örneğin bir küme,
NP içindeyse ve her NP problemi aşağıda betimlenene benzer, fakat ek olarak
polinom zamanda hesaplanabilen bir indirgeme kullanılarak ona
indirgenebiliyorsa ``NP-tam''dır.

İndirgeme kavramını örtük olarak zaten kullandık. $K$ kümesini
\[
  K=\Setabs{x}{\cfind{x}(x)\fdefined},
\]
yani $K=\Setabs{x}{x\in W_x}$ olarak tanımlayalım. Durma probleminin
çözülemez olduğuna ilişkin ispatımız (\olref[hlt]{thm:halting-problem}) en
doğrudan biçimde $K$'nin hesaplanabilir olmadığını gösterir. $K_0$'ın
\[
  K_0=\Setabs{\tuple{e,x}}{\cfind{e}(x)\fdefined},
\]
yani $K_0=\Setabs{\tuple{e,x}}{x\in W_e}$ olduğunu anımsayın. $K$'nin
hesaplanamazlığına ilişkin her ispatı $K_0$'ın hesaplanamazlığına genişletmek
kolaydır: $K_0$ hesaplanabilir olsaydı !!a{element}~$x$'in $K$ içinde bulunup
bulunmadığına yalnızca $\tuple{x,x}$ çiftinin $K_0$ içinde bulunup
bulunmadığını sorarak karar verebilirdik. $x$'i $\tuple{x,x}$ çiftine
götüren $f$ fonksiyonu, $K$'nin $K_0$'a bir \emph{indirgenmesi} örneğidir.
\end{explain}

\begin{defn}
$A$ ile $B$ doğal sayı kümeleri olsun. Hesaplanabilir
$f\colon\Nat\to\Nat$ fonksiyonu, her doğal sayı $x$ için
\[
  x\in A \quad\text{ancak ve ancak}\quad f(x)\in B
\]
olduğunda $A$'nın $B$'ye \emph{çok-bir indirgenmesi}dir. Böyle bir $f$
indirgemesi varsa $A$, $B$'ye \emph{çok-bir indirgenebilir}dir ve
$A\leq_m B$ yazarız. $A$, $B$'ye ve $B$, $A$'ya çok-bir indirgenebilirse
$A$ ile $B$ \emph{çok-bir eşdeğer}dir; bunu $A\equiv_m B$ ile gösteririz.
\end{defn}

Yukarıdaki tanımdaki $f$ fonksiyonu ayrıca bire birse $A$'ya $B$'ye
\emph{bir-bir indirgenebilir} denir. Aşağıda betimlenen indirgemelerin çoğu
bu daha güçlü koşulu karşılar; fakat bu olguyu kullanmayacağız.

\begin{digress}
Bir-bir indirgenebilirliğin gerçekten çok-bir indirgenebilirlikten daha güçlü
bir koşul olması doğru, fakat hiç de açık değildir. Başka bir deyişle,
$A$'nın $B$'ye çok-bir indirgenebilir fakat bir-bir indirgenemez olduğu sonsuz
$A$ ve $B$ kümeleri vardır.
\end{digress}

